% ---------------------------------------------------------------------
% 2.7 Procesamiento digital de señales biomédicas
% ---------------------------------------------------------------------
\section[PROCESAMIENTO DIGITAL DE SEÑALES BIOMÉDICAS]{Procesamiento Digital de Señales
Biomédicas}
\label{sec:pds}

El preprocesamiento de las señales EEG constituye una etapa determinante para el desempeño de
los modelos de clasificación, dado que estas señales presentan una relación señal-ruido baja y
son susceptibles a artefactos de origen fisiológico y ambiental. Para atender esta problemática
se emplean técnicas de filtrado pasa banda, orientadas a delimitar el rango de frecuencias
asociado a la componente P300, así como la segmentación de la señal continua en ventanas
temporales alrededor de cada estímulo y la normalización de las mismas, cuyos fundamentos se
describen a lo largo de esta sección.

\subsection{Muestreo y teorema de Nyquist}

Las señales de origen fisiológico son de naturaleza analógica, es decir, están definidas en todo
instante de tiempo y su amplitud puede tomar cualquier valor dentro de su rango dinámico; para
procesarlas en una computadora es necesario muestrearlas a intervalos regulares
$T_s = 1/f_s$, obteniendo la secuencia discreta $x(n) = x_a(nT_s)$ \cite{proakis2007}. El teorema
de muestreo establece que una señal cuyo contenido en frecuencia está limitado a $f_{\max}$ puede
reconstruirse sin pérdida a partir de sus muestras siempre que
\begin{equation}
f_s \geq 2 f_{\max},
\label{eq:nyquist}
\end{equation}
y si esta condición no se cumple, las componentes de frecuencia superiores a $f_s/2$, conocida
como frecuencia de Nyquist, se reflejan sobre las frecuencias bajas en un fenómeno llamado
\textit{aliasing} \cite{oppenheim2010}. La base de datos utilizada en este trabajo fue muestreada a
240~Hz, por lo que su frecuencia de Nyquist de 120~Hz está muy por encima de la banda de interés
de la P300, que se concentra por debajo de 20~Hz.

\subsection{Transformada de Fourier y densidad espectral de potencia}

La transformada discreta de Fourier (DFT) permite representar una secuencia de $N$ muestras como
una suma de componentes sinusoidales, y se define como
\begin{equation}
X(k) = \sum_{n=0}^{N-1} x(n)\, e^{-j 2\pi k n / N}, \qquad k = 0, 1, \dots, N-1,
\label{eq:dft}
\end{equation}
donde cada coeficiente $X(k)$ corresponde a la frecuencia $f_k = k f_s / N$; su cálculo directo
requiere del orden de $N^2$ operaciones, mientras que el algoritmo de la transformada rápida de
Fourier (FFT) lo reduce a $N\log_2 N$ \cite{proakis2007}. A partir de ella se estima la densidad
espectral de potencia (PSD), que describe cómo se distribuye la potencia de la señal en función
de la frecuencia. En el caso del EEG suele emplearse el método de Welch, que divide la señal en
segmentos traslapados, calcula el periodograma de cada uno y los promedia, reduciendo así la
varianza de la estimación a costa de una menor resolución en frecuencia \cite{oppenheim2010}.
Contreras utilizó este análisis para identificar que la información útil para detectar la P300 en
sus registros se encontraba entre 1 y 25~Hz \cite{contreras2021}; la Figura~\ref{fig:psd} muestra
una PSD típica en la que se distinguen el pico del ritmo alfa y la interferencia de la red
eléctrica.

\begin{figure}[htbp]
\centering
\includegraphics[width=0.8\textwidth]{02-marco-referencial/psd_welch.pdf}
\caption{Densidad espectral de potencia estimada con el método de Welch para una señal de EEG
sintética. Elaboración propia.}
\label{fig:psd}
\end{figure}

\subsection{Convolución y sistemas lineales}

Un sistema lineal e invariante en el tiempo queda completamente descrito por su respuesta al
impulso $h(n)$, y su salida ante cualquier entrada $x(n)$ se obtiene mediante la operación de
convolución \cite{proakis2007}:
\begin{equation}
y(n) = (x * h)(n) = \sum_{k=-\infty}^{\infty} x(k)\, h(n-k).
\label{eq:convolucion}
\end{equation}
Esta operación es la base de los filtros digitales que se describen a continuación y, como se
verá en la Sección~\ref{sec:cnn}, también de las capas de las redes neuronales convolucionales,
en las que los coeficientes de $h(n)$ ya no se diseñan manualmente sino que se aprenden a partir
de los datos \cite{goodfellow2016}.

\subsection{Filtros digitales}

Los filtros digitales se clasifican en dos grandes familias según la duración de su respuesta al
impulso. Los filtros de respuesta finita al impulso (FIR) calculan la salida únicamente a partir de
las muestras de entrada, son siempre estables y pueden diseñarse con fase lineal, mientras que los
de respuesta infinita al impulso (IIR) utilizan además salidas anteriores, lo que les permite
lograr una transición más abrupta con un orden mucho menor, aunque con una respuesta de fase no
lineal \cite{proakis2007,oppenheim2010}. Entre los filtros IIR, el de Butterworth es uno de los más
utilizados en EEG porque su respuesta en magnitud es máximamente plana en la banda de paso
\cite{oppenheim2010}. De acuerdo con la banda que dejan pasar, los filtros pueden ser pasa bajas,
pasa altas, pasa banda o rechaza banda; un caso particular de estos últimos es el filtro
\textit{notch}, diseñado para eliminar una sola frecuencia como la de la red eléctrica. La
Tabla~\ref{tab:filtros} resume sus características y la Figura~\ref{fig:filtros} muestra la
respuesta en frecuencia de un filtro pasa banda entre 0.1 y 20~Hz, como el aplicado en el trabajo
preliminar, y de un filtro \textit{notch} en 60~Hz.

\begin{table}[htbp]
\centering
\caption{Comparación de filtros digitales FIR e IIR}
\label{tab:filtros}
\small
\begin{tabularx}{\textwidth}{L{3.3cm}YY}
\toprule
\textbf{Característica} & \textbf{FIR} & \textbf{IIR (p.\,ej. Butterworth)} \\
\midrule
Ecuación en diferencias & $y(n)=\sum_{k=0}^{M} b_k\,x(n-k)$ & $y(n)=\sum_{k=0}^{M} b_k\,x(n-k)-\sum_{k=1}^{N} a_k\,y(n-k)$ \\
Estabilidad & Siempre estable. & Depende de la ubicación de los polos. \\
Fase & Puede ser lineal (retardo constante). & No lineal; se compensa con filtrado hacia adelante y hacia atrás. \\
Orden necesario & Alto para transiciones abruptas. & Bajo para la misma selectividad. \\
Costo computacional & Mayor. & Menor. \\
\bottomrule
\end{tabularx}
\par\smallskip{\footnotesize Elaboración propia con información de \cite{proakis2007,oppenheim2010}.}
\end{table}

\begin{figure}[htbp]
\centering
\includegraphics[width=\textwidth]{02-marco-referencial/respuesta_filtros.pdf}
\caption{Respuesta en magnitud de (a) un filtro pasa banda Butterworth de cuarto orden entre 0.1 y
20~Hz y (b) un filtro \textit{notch} en 60~Hz, para $f_s = 240$~Hz. Elaboración propia.}
\label{fig:filtros}
\end{figure}

\subsection{Filtrado de fase cero y remuestreo}

La fase no lineal de los filtros IIR provoca que distintas frecuencias se retrasen de forma
diferente, lo que puede desplazar el pico de un ERP y alterar la latencia medida de la P300
\cite{luck2014}. Cuando el procesamiento se realiza fuera de línea, este efecto se elimina
aplicando el filtro dos veces, una en sentido directo y otra en sentido inverso, técnica conocida
como filtrado de fase cero, que duplica la atenuación efectiva del filtro sin introducir retardo
\cite{oppenheim2010}. Por otra parte, el remuestreo permite reducir la frecuencia de muestreo una
vez que la señal ya fue filtrada; en la decimación por un factor $M$ se conserva una de cada $M$
muestras, siempre que antes se haya aplicado un filtro antialiasing que elimine el contenido por
encima de la nueva frecuencia de Nyquist \cite{proakis2007}. En el trabajo preliminar se redujo la
frecuencia de 240~Hz a 120~Hz, lo que disminuye a la mitad el tamaño de cada época sin perder
información en la banda de 0.1 a 20~Hz.

\subsection{Referencia y corrección de línea base}

Además del filtrado en frecuencia, es común modificar la referencia de los canales después del
registro. La referencia promedio común resta a cada canal $x_{ch}(n)$ el promedio de los $C$
canales en cada instante,
\begin{equation}
x^{\text{CAR}}_{ch}(n) = x_{ch}(n) - \frac{1}{C}\sum_{c=1}^{C} x_{c}(n),
\label{eq:car}
\end{equation}
con lo que se atenúan las componentes comunes a todos los electrodos, como las derivadas de la
propia referencia \cite{luck2014}. De forma complementaria, la corrección de línea base resta a
cada época el valor promedio de un intervalo previo al estímulo, típicamente de 100 a 200~ms, para
que las amplitudes de los componentes del ERP se midan respecto a un nivel común \cite{luck2014}.

\subsection{Segmentación en épocas}

La segmentación consiste en dividir la señal continua en ventanas sincronizadas con cada estímulo,
que constituyen las muestras de entrada al clasificador. La longitud de la ventana debe ser
suficiente para contener la P300 completa; Contreras empleó ventanas de 1005~ms a partir del
estímulo \cite{contreras2021}, mientras que en el trabajo preliminar se utilizaron ventanas de 78
muestras, equivalentes a unos 650~ms a 120~Hz. Dado que el intervalo entre inicios de estímulos es
de solo 175~ms, las épocas de intensificaciones consecutivas se traslapan, como se observa en la
Figura~\ref{fig:epocas}, de modo que una época no objetivo puede contener parte de la P300
provocada por una intensificación objetivo cercana.

\begin{figure}[htbp]
\centering
\includegraphics[width=\textwidth]{02-marco-referencial/segmentacion_epocas.pdf}
\caption{Segmentación de la señal continua en épocas de 0 a 650~ms a partir de cada
intensificación. Señal sintética, elaboración propia.}
\label{fig:epocas}
\end{figure}

\subsection{Normalización}

Las amplitudes del EEG varían entre canales, sesiones y sujetos, por lo que antes de entrenar un
modelo se acostumbra normalizar los datos. La normalización \textit{z-score} transforma cada canal
para que tenga media cero y desviación estándar unitaria, mientras que la normalización
\textit{min-max}, utilizada por Contreras \cite{contreras2021}, escala los valores al intervalo
$[0,1]$:
\begin{equation}
z(n) = \frac{x(n)-\mu}{\sigma}, \qquad
x'(n) = \frac{x(n)-\min(x)}{\max(x)-\min(x)}.
\label{eq:normalizacion}
\end{equation}
La primera es menos sensible a valores atípicos, lo cual resulta conveniente en presencia de
artefactos de gran amplitud, y fue la empleada de manera independiente por canal en el trabajo
preliminar para evitar que los electrodos con mayor amplitud dominen el aprendizaje de la red.

\subsection{Remoción de artefactos}

Cuando los artefactos se traslapan en frecuencia con la señal de interés, como ocurre con los
parpadeos y la banda de la P300, el filtrado no basta para eliminarlos. La estrategia más simple
consiste en descartar las épocas cuya amplitud supere un umbral, aunque esto reduce la cantidad de
datos disponibles; como alternativa, el análisis de componentes independientes (ICA) separa la
señal multicanal en fuentes estadísticamente independientes, de modo que las asociadas a
movimientos oculares o musculares pueden identificarse y eliminarse antes de reconstruir el EEG
\cite{jung2000}. Esta técnica se popularizó gracias a herramientas como EEGLAB, que la integran
junto con otras rutinas de análisis de ensayos individuales \cite{delorme2004}.

\subsection{Filtrado espacial}

Finalmente, además de filtrar cada canal en el tiempo, es posible combinar los canales mediante un
filtro espacial que realce la señal de interés. El algoritmo xDAWN, propuesto por Rivet et al.,
estima los filtros espaciales que maximizan la relación entre la señal evocada y la señal más
ruido, y ha mostrado mejorar la detección de la P300 aun con un número reducido de canales
\cite{rivet2009}. Este tipo de filtrado es, en cierto sentido, lo que las capas convolucionales
espaciales de redes como EEGNet aprenden de manera automática a partir de los datos
\cite{lawhern2018}.
